% R. Brown's translations of two letters to L. Breen, with his note and comments (cote 134-2): archive cote 134-2, site page 138 (PDF page 139).
% translation of the letter of 17-19 July 1975. Draft compared: ps1.tex, lines 1585-1612.
% AI-transcribed by Claude Opus 5.5 (claude-opus-5-5, adaptive thinking, effort high), Message Batches API, batch brown_full (msgbatch_01Md3i5yK1PxkpJn5GHfEExE), 2026-11-13; not yet checked by eye.
\scanpage{134-2/138}
\mspage{- 11 -}

where ``literally'' $p''$ is now an inclusion, but ``morally'', it is
a \uline{fibration} with a very pretty fibre [notably in finite dimensions,
these are non-singular varieties - this is then much better than
that which is given by the yoga of Cartan-Serre ``every continuous
mapping is equivalent to a fibration''...] . This last diagram
is acceptable to a direct construction, purely algebraic, in the
context of schemes, once one has developed the construction of
étale tubular neighbourhood.

The point I wish to come to, is that the consideration of
sheaves (or complexes of stacks, of n-stacks ...) which are $P$-%
constructible on an $X$, where $P$ is a given ``equi singular''
stratification, comes in actual particular cases to the knowledge
of a diagram of ordinary \uline{homotopy} \uline{types} (or pro-types, if one
comes back to the étale topology)
\diagram{\begin{tikzcd}
  \Delta \ar[r,"j"]\ar[d,swap,"p"] & X_1 \\
  X_0 &
\end{tikzcd}}
and taking local coefficients systems (or locally constant n-stacks)
on the vertices $X_0$ , $X_1$ , related by homeomorphisms of compat-%
ibility of the type $p^*(F_0) \to j^*(F_1)$ . It should be an amusing
excercise (which I have not yet done) to verify and to make
explicit how the usual operations on sheaves (on $X$ , or on a union
of strata of $X$ ) can be expressed in this dictionary, in the
stated case of non-singular strata (at this moment, there will be
a difficulty with the dualising complexes,

\mspage{\hand{139}}
